\section{Transmission-Line Theory Behind TDR}
\label{sec:theory}

\subsection{Telegrapher's Equations and Lossy Wave Propagation}
A coaxial cable supports a transverse electromagnetic (TEM) wave. Along a uniform line, the voltage and current phasors obey the telegrapher's equations \cite{pozar2011microwave}:
\begin{equation}
\frac{d V(z, \omega)}{dz} = -Z(\omega) I(z, \omega), \quad \frac{d I(z, \omega)}{dz} = -Y(\omega) V(z, \omega)
\label{eq:telegrapher}
\end{equation}
where $Z(\omega) = R(\omega) + j\omega L$ and $Y(\omega) = G(\omega) + j\omega C$ represent the series impedance and shunt admittance per unit length, respectively. Eliminating $I$ gives the wave equation $d^2V/dz^2 = \gamma^2 V$, whose solutions are forward- and backward-traveling waves with complex propagation constant $\gamma(\omega) = \alpha(\omega) + j\beta(\omega) = \sqrt{Z(\omega)Y(\omega)}$ and characteristic impedance $Z_0(\omega) = \sqrt{Z(\omega)/Y(\omega)}$:
\begin{equation}
V(z, \omega) = V_0^+ e^{-\gamma(\omega) z} + V_0^- e^{+\gamma(\omega) z}
\label{eq:wave_sol}
\end{equation}

\subsection{Lossy Coaxial Cable Physics (Skin Effect \& Dielectric Dissipation)}
A real coaxial cable loses energy in both its conductors and its dielectric, and both losses grow with frequency. We model them as follows:
\begin{enumerate}
    \item \textbf{Skin Effect:} At radio frequencies, current concentrates within a surface layer of depth $\delta_s = \sqrt{2/(\omega \mu \sigma)}$. Consequently, the series resistance increases with the square root of frequency:
    \begin{equation}
    R(\omega) = R_{\mathrm{dc}} + k_{\mathrm{skin}} \sqrt{\frac{\omega}{2\pi}}
    \label{eq:skin_effect}
    \end{equation}
    where $R_{\mathrm{dc}} = 0.04\ \Omega/\mathrm{m}$ is the DC resistance and $k_{\mathrm{skin}} = 1.5 \times 10^{-4}\ \Omega/(\mathrm{m}\cdot\sqrt{\mathrm{Hz}})$ sets the strength of the skin effect (about 0.13~dB/m at 100~MHz).
    \item \textbf{Dielectric Loss:} The insulating polyethylene dielectric possesses a non-zero loss tangent $\tan\delta$, yielding a shunt conductance directly proportional to frequency:
    \begin{equation}
    G(\omega) = \omega C \tan\delta
    \label{eq:dielectric_loss}
    \end{equation}
    with $\tan\delta = 2.0 \times 10^{-4}$.
\end{enumerate}
The nominal phase velocity in the TEM mode is $u = v_f c = 1/\sqrt{LC}$, where $v_f \approx 0.66$ denotes the cable velocity factor and $c = 2.9979 \times 10^8\ \mathrm{m/s}$.

\subsection{Reflection Coefficients and Discontinuity Boundary Conditions}
When a propagating voltage wave encounters a localized impedance boundary $Z_{\mathrm{load}}$ at position $z$, the continuity of voltage and current dictates a reflected wave with voltage reflection coefficient $\Gamma$:
\begin{equation}
\Gamma = \frac{Z_{\mathrm{load}} - Z_0}{Z_{\mathrm{load}} + Z_0}
\label{eq:gamma}
\end{equation}

\textbf{Hard faults.} A short ($Z_f \to 0$) gives $\Gamma = -1$: the echo has equal size and opposite sign, and it pulls the sending-end voltage to zero when it arrives. An open ($Z_f \to \infty$) gives $\Gamma = +1$ and doubles the voltage.

\textbf{Soft faults.}
\begin{itemize}
    \item \textbf{Shunt Faults ($Z_f$ in parallel with line):} The effective impedance seen at the discontinuity is the parallel combination of $Z_f$ and the downstream line impedance $Z_0$:
    \begin{equation}
    Z_{\mathrm{eff}} = \frac{Z_f Z_0}{Z_f + Z_0} \implies \Gamma_{\mathrm{shunt}} = \frac{-Z_0}{2 Z_f + Z_0}
    \label{eq:gamma_shunt}
    \end{equation}
    For a shunt resistor $R_f$, $\Gamma_{\mathrm{shunt}} \in (-1, 0)$ produces a negative step echo. For a shunt capacitor $C_f$, the uncharged capacitor acts as an instantaneous short ($\Gamma \to -1$), recovering exponentially to $\Gamma \to 0$ as it charges, yielding a characteristic transient negative spike.
    \item \textbf{Series Faults ($R_f$ in series with inner conductor):} The downstream impedance is $Z_0 + R_f$:
    \begin{equation}
    \Gamma_{\mathrm{series}} = \frac{(Z_0 + R_f) - Z_0}{(Z_0 + R_f) + Z_0} = \frac{R_f}{2 Z_0 + R_f}
    \label{eq:gamma_series}
    \end{equation}
    Here $\Gamma_{\mathrm{series}} \in (0, 1)$, generating a positive step echo.
    \item \textbf{Characteristic Impedance Section ($z_0\_section$):} A finite segment of length $w$ with altered geometry has local impedance $Z_0'$. Reflections occur at both the entry ($z=d$) with $\Gamma_1 = (Z_0' - Z_0)/(Z_0' + Z_0)$ and exit ($z=d+w$) with $\Gamma_2 = -\Gamma_1$, creating a rectangular pulse or flat pedestal of duration $\Delta t = 2w/u$.
\end{itemize}

\subsection{Complex Load Terminations and Reactive Time Constants}
At the terminal end of the line ($z=l$), we evaluate six load classes:
\begin{enumerate}
    \item \textbf{Matched Load ($Z_L = Z_0$):} $\Gamma_L = 0$; no reflection returns from the load.
    \item \textbf{Open and Short ($Z_L \to \infty$, $Z_L = 0$):} $\Gamma_L = +1$ and $-1$.
    \item \textbf{Resistive Load ($Z_L = R_L$):} $\Gamma_L = (R_L - Z_0)/(R_L + Z_0)$.
    \item \textbf{Inductive-Resistive Load ($RL$ Series):} The load impedance $Z_L(s) = R_L + sL$ yields the reflection coefficient in the Laplace domain:
    \begin{equation}
    \Gamma_L(s) = \frac{sL + R_L - Z_0}{sL + R_L + Z_0} = 1 - \frac{2 Z_0 / L}{s + (R_L + Z_0)/L}
    \end{equation}
    In response to an incident step $V_{\mathrm{inc}}$, the reflected waveform exhibits an initial open-circuit-like step ($\Gamma \to +1$) decaying exponentially toward the resistive steady-state $\Gamma_\infty = (R_L - Z_0)/(R_L + Z_0)$ with time constant:
    \begin{equation}
    \tau_{\mathrm{RL}} = \frac{L}{R_L + Z_0}
    \label{eq:tau_rl}
    \end{equation}
    \item \textbf{Capacitive-Resistive Load ($RC$ Parallel):} A resistor in parallel with a capacitor, $Z_L(s) = R_L/(1 + sR_LC)$, gives
    \begin{equation}
    \Gamma_L(s) = \frac{(R_L - Z_0) - s C R_L Z_0}{(R_L + Z_0) + s C R_L Z_0}
    \end{equation}
    At the instant the step arrives the uncharged capacitor acts as a short ($\Gamma_0 = -1$), so the trace dips; as it charges the reflection relaxes to the resistive value $\Gamma_\infty = (R_L - Z_0)/(R_L + Z_0)$ with the time constant of $C$ discharging through $R_L \parallel Z_0$:
    \begin{equation}
    \tau_{\mathrm{RC}} = C \, \frac{R_L Z_0}{R_L + Z_0}
    \label{eq:tau_rc}
    \end{equation}
\end{enumerate}

\subsection{Lattice Bounce Diagram \& Physics Masks}
In a finite line of length $l$ fed by a generator $V_g(t)$ with source resistance $R_g$, the initial incident wave launched into the line at $t=0$ is:
\begin{equation}
V_{\mathrm{inc}} = V_g \cdot \frac{Z_0}{Z_0 + R_g}
\label{eq:v_inc}
\end{equation}
The wave propagates toward the fault at $d$ and load at $l$, undergoing successive partial reflections and transmissions at the discontinuity, load ($\Gamma_L$), and generator ($\Gamma_g = (R_g - Z_0)/(R_g + Z_0)$).
The round-trip transit time is defined as:
\begin{equation}
T_{\mathrm{nom}} = \frac{2 l}{u} = \frac{2 l}{v_f c}
\label{eq:t_nom}
\end{equation}
A fault at distance $d$ produces its first echo at $t_f = 2d/u = (d/l) T_{\mathrm{nom}}$, while the load echo arrives at $t_L = T_{\mathrm{nom}}$.

\textbf{Consequences for the learning tasks.} Two targets do not always exist physically. On a fault-free line there is no fault position, so location is neither trained nor scored for \texttt{none}. A hard short or open reflects the whole wave, so nothing reaches the load and it cannot be observed. The load tasks are therefore masked out for \texttt{short} and \texttt{open} faults, both in the training loss and in the metrics.

